\section{Accessing the fields}\label{sec:06-groups:05-accessing-fields:1}

Because \texttt{intGroup} is a term of type \texttt{Group Int}, its fields can be projected out exactly as in Chapter 2:

\begin{lstlisting}
#eval intGroup.op 3 4        -- 7
#eval intGroup.id             -- 0
#eval intGroup.inv 5          -- -5

#check intGroup.assoc         -- a proof, for every a b c, of associativity
\end{lstlisting}

\begin{mathreading}

The projections recover the individual components of the structure: \texttt{intGroup.op} is the multiplication \(\cdot\) (so \texttt{intGroup.op 3 4} is \(3 + 4 = 7\) in \(\mathbb{Z}\)), \texttt{intGroup.id} is \(e = 0\), and \texttt{intGroup.inv} is \((-)^{-1} = -(-)\). The key point is that \texttt{intGroup.assoc} projects out a \emph{proof}: it is the element of \(\forall
a,b,c,\ (a\cdot b)\cdot c = a\cdot(b\cdot c)\) that was supplied when building the group. Data-fields and proof-fields are accessed the same way because, in the dependent-pair view (a \texttt{structure} is, underneath, exactly this kind of dependent pair), both are just coordinates of the same tuple.

\end{mathreading}

\textbf{Mathlib equivalent.} There is no \texttt{intGroup.op 3 4}-style field access to write at all. Once \texttt{Int} is known to be an \href{https://loogle.lean-lang.org/?q=AddCommGroup}{\texttt{AddCommGroup}}, the ordinary \texttt{+}/\texttt{0}/\texttt{-\allowbreak{}} notations already resolve to that instance's operations directly:

\begin{lstlisting}
#eval (3 : Int) + 4
#eval (0 : Int)
#eval -(5 : Int)
#check (add_assoc : ∀ a b c : Int, (a + b) + c = a + (b + c))
\end{lstlisting}

This is the same contrast as §3: the book's \texttt{intGroup.op}/\texttt{.id}/\texttt{.inv} are projections out of a bundle built by hand, while Mathlib's \texttt{+}/\texttt{0}/ \texttt{-\allowbreak{}} are notation that the typeclass system has already wired to the right instance. The underlying ``which \texttt{AddCommGroup} instance is this?'' bookkeeping remains invisible unless sought out (for example, with \texttt{\#print}).
